% Sections won't have numbers
% Make the page area as large as possible
% Set the first level numbering to arabic and the second level
% to alphabetic. The remaining two levels are arabic
%Change to \pagestyle{empty} to suppress numbers on following pages
%\pagestyle{empty}
% Do not reset outermost enumerate counter


\documentclass[onecolumn]{article}
%%%%%%%%%%%%%%%%%%%%%%%%%%%%%%%%%%%%%%%%%%%%%%%%%%%%%%%%%%%%%%%%%%%%%%%%%%%%%%%%%%%%%%%%%%%%%%%%%%%%%%%%%%%%%%%%%%%%%%%%%%%%%%%%%%%%%%%%%%%%%%%%%%%%%%%%%%%%%%%%%%%%%%%%%%%%%%%%%%%%%%%%%%%%%%%%%%%%%%%%%%%%%%%%%%%%%%%%%%%%%%%%%%%%%%%%%%%%%%%%%%%%%%%%%%%%
\usepackage{amsmath}
\usepackage{sw20exm2}
\usepackage[compat2]{geometry}

\setcounter{MaxMatrixCols}{10}
%TCIDATA{OutputFilter=LATEX.DLL}
%TCIDATA{Version=5.50.0.2953}
%TCIDATA{<META NAME="SaveForMode" CONTENT="1">}
%TCIDATA{BibliographyScheme=Manual}
%TCIDATA{Created=Sunday, August 26, 2012 12:26:36}
%TCIDATA{LastRevised=Monday, June 08, 2026 10:47:53}
%TCIDATA{<META NAME="ViewSettings" CONTENT="0">}
%TCIDATA{<META NAME="GraphicsSave" CONTENT="32">}
%TCIDATA{<META NAME="DocumentShell" CONTENT="Exams and Syllabi\SW\SW Exam #1 - 8.5x14 Page">}
%TCIDATA{CSTFile=Exam.cst}

\setlength{\topmargin}{-0.75in}
\setlength{\textheight}{12.25in}
\setlength{\oddsidemargin}{0.0in}
\setlength{\evensidemargin}{0.0in}
\setlength{\textwidth}{6.5in}
\def\labelenumi{\arabic{enumi}.}
\def\theenumi{\arabic{enumi}}
\def\labelenumii{(\alph{enumii})}
\def\theenumii{\alph{enumii}}
\def\p@enumii{\theenumi.}
\def\labelenumiii{\arabic{enumiii}.}
\def\theenumiii{\arabic{enumiii}}
\def\p@enumiii{(\theenumi)(\theenumii)}
\def\labelenumiv{\arabic{enumiv}.}
\def\theenumiv{\arabic{enumiv}}
\def\p@enumiv{\p@enumiii.\theenumiii}
\pagestyle{plain}
\setcounter{secnumdepth}{0}
\newtheorem{theorem}{Theorem}
\newtheorem{acknowledgement}[theorem]{Acknowledgement}
\newtheorem{algorithm}[theorem]{Algorithm}
\newtheorem{axiom}[theorem]{Axiom}
\newtheorem{case}[theorem]{Case}
\newtheorem{claim}[theorem]{Claim}
\newtheorem{conclusion}[theorem]{Conclusion}
\newtheorem{condition}[theorem]{Condition}
\newtheorem{conjecture}[theorem]{Conjecture}
\newtheorem{corollary}[theorem]{Corollary}
\newtheorem{criterion}[theorem]{Criterion}
\newtheorem{definition}[theorem]{Definition}
\newtheorem{example}[theorem]{Example}
\newtheorem{exercise}[theorem]{Exercise}
\newtheorem{lemma}[theorem]{Lemma}
\newtheorem{notation}[theorem]{Notation}
\newtheorem{problem}[theorem]{Problem}
\newtheorem{proposition}[theorem]{Proposition}
\newtheorem{remark}[theorem]{Remark}
\newtheorem{solution}[theorem]{Solution}
\newtheorem{summary}[theorem]{Summary}
\newenvironment{proof}[1][Proof]{\noindent\textbf{#1.} }{\ \rule{0.5em}{0.5em}}
\input{tcilatex}
\begin{document}


\begin{center}
\begin{equation*}
\FRAME{itbpF}{1.0525in}{0.7515in}{0in}{}{}{Figure}{\special{language
"Scientific Word";type "GRAPHIC";maintain-aspect-ratio TRUE;display
"USEDEF";valid_file "T";width 1.0525in;height 0.7515in;depth
0in;original-width 2.3168in;original-height 1.6466in;cropleft "0";croptop
"1";cropright "1";cropbottom "0";tempfilename
'TCTQIQ04.bmp';tempfile-properties "XPR";}}
\end{equation*}%
\textbf{F\'{\i}sica II (2026-1)}

\textbf{Tarea 3}\emph{: Condensadores, Leyes de Kirchhoff y Fuerza magn\'{e}%
tica}

\textit{Profesor: Arturo Fern\'{a}ndez P\'{e}rez}

\textbf{Plazo de Entrega: S\'{a}bado 13 de junio a las 20:00 hrs.}
\end{center}

\begin{enumerate}
\item (2 puntos) En el circuito de la Fig. 1, $C_{1}=12,0$ [pF], $%
C_{2}=C_{5}=5,0$ $[$pF$]$ y $C_{6}=C_{7}=8,0$ [pF]. El capacitor $C_{3}$ es
un condensador de placas planas \textit{de forma cuadrada}, con lado $l=5,0$
[cm] y distancia entre placas $d=2,0$ [mm]. Entre sus placas hay 3 diel\'{e}%
ctricos con constantes $\kappa _{1}=2,$ $\kappa _{2}=3,$ $\kappa _{3}=5$, en
la distribuci\'{o}n que se muestra en la Fig. 1. Por otro lado, el capacitor 
$C_{4}$ es un condensador cil\'{\i}ndrico de largo $\ell =5,0$ [cm], radio
interior $r_{a}=2,0$ [mm] y radio exterior $r_{b}=3,0$ [mm], con un diel\'{e}%
ctrico de constante $\kappa _{1}=2\,\ $entre sus electrodos. Si $V_{ab}=12$
[V],

\begin{enumerate}
\item Determine las capacitancias $C_{3}$ y $C_{4}.$

\item Determine la capacitancia equivalente $C_{eq}$ del circuito.

\item Determine la carga $q$ y la energ\'{\i}a almacenada $U$ en cada
condensador. 
\begin{equation*}
\FRAME{itbpFU}{6.0494in}{2.0946in}{0.1254in}{\Qcb{Figura 1}}{}{Figure}{%
\special{language "Scientific Word";type "GRAPHIC";maintain-aspect-ratio
TRUE;display "USEDEF";valid_file "T";width 6.0494in;height 2.0946in;depth
0.1254in;original-width 14.976in;original-height 5.1586in;cropleft
"0";croptop "1";cropright "1";cropbottom "0";tempfilename
'TGBI6E01.bmp';tempfile-properties "XPR";}}
\end{equation*}
\end{enumerate}

\item (1 punto) Considerando el circuito de la Fig. 2, utilize las leyes
Kirchhoff para completar la Tabla 1.%
\begin{equation*}
\FRAME{itbpFU}{3.4912in}{0.8743in}{0in}{\Qcb{Tabla 1}}{}{Figure}{\special%
{language "Scientific Word";type "GRAPHIC";maintain-aspect-ratio
TRUE;display "USEDEF";valid_file "T";width 3.4912in;height 0.8743in;depth
0in;original-width 6.6565in;original-height 1.6457in;cropleft "0";croptop
"1";cropright "1";cropbottom "0";tempfilename
'TGA9GN04.wmf';tempfile-properties "XPR";}}
\end{equation*}

\begin{equation*}
\FRAME{itbpFU}{2.2122in}{0.9868in}{0in}{\Qcb{Figura 2}}{}{Figure}{\special%
{language "Scientific Word";type "GRAPHIC";maintain-aspect-ratio
TRUE;display "USEDEF";valid_file "T";width 2.2122in;height 0.9868in;depth
0in;original-width 6.3858in;original-height 2.8331in;cropleft "0";croptop
"1";cropright "1";cropbottom "0";tempfilename
'TGA9GB03.wmf';tempfile-properties "XPR";}}
\end{equation*}%
\pagebreak 

\item (1 punto) La Fig. 3 muestra un circuito que combina resistencias y
bater\'{\i}as en distintas configuraciones. Considere los sentidos de las
corrientes $I_{1}$, $I_{2}$ e $I_{3}$, as\'{\i} como el recorrido de c/u \
de las mallas, indicados como $M_{1}$ y $M_{2}$.

\begin{enumerate}
\item Simplifique el circuito todo lo posible obteniendo las resistencias
equivalentes $R_{eq}$ y dibuje el esquema reducido resultante.

\item Utilizando las leyes de Kirchhoff, determine los valores de $I_{1}$, $%
I_{2}$ e $I_{3}$.%
\begin{equation*}
\FRAME{itbpFU}{3.5604in}{1.9285in}{0in}{\Qcb{Figura 3}}{}{Figure}{\special%
{language "Scientific Word";type "GRAPHIC";maintain-aspect-ratio
TRUE;display "USEDEF";valid_file "T";width 3.5604in;height 1.9285in;depth
0in;original-width 7.3129in;original-height 3.9479in;cropleft "0";croptop
"1";cropright "1";cropbottom "0";tempfilename
'TGA9MA05.wmf';tempfile-properties "XPR";}}
\end{equation*}
\end{enumerate}

\item (1 punto) La Fig. 4 representa un lazo de corriente cuadrado de lado $%
a=10,0$ [cm] por donde circula una corriente $I_{0}=1,8$ [A]. Los sentidos
de circulaci\'{o}n de las corrientes en cada lado del cuadrado est\'{a}n
indicadas en la Fig. 4, as\'{\i} como el sistema de referencia. Si este lazo
de corriente est\'{a} sujeto a un campo magn\'{e}tico $\vec{B}=1,2\hat{\jmath%
}+0.7\hat{k}$ [T], determine:

\begin{enumerate}
\item La fuerza magn\'{e}tica $\vec{F}_{m}$ sobre c/u de los lados del
cuadrado.

\item El momento magn\'{e}tico $\vec{\mu}$ del lazo de corriente.

\item El torque magn\'{e}tico $\vec{\tau}_{m}$ inducido por el campo magn%
\'{e}tico $\vec{B}$ sobre el lazo de corriente. 
\begin{equation*}
\FRAME{itbpFU}{3.5492in}{2.0124in}{0in}{\Qcb{Figura 4}}{}{Figure}{\special%
{language "Scientific Word";type "GRAPHIC";maintain-aspect-ratio
TRUE;display "USEDEF";valid_file "T";width 3.5492in;height 2.0124in;depth
0in;original-width 6.0105in;original-height 3.3961in;cropleft "0";croptop
"1";cropright "1";cropbottom "0";tempfilename
'TGA9W707.wmf';tempfile-properties "XPR";}}
\end{equation*}
\end{enumerate}
\end{enumerate}

\end{document}
